
\section{Efectividad y configuración}
Una tarea puede depender de modelo, serie, modificaciones y componentes. Un párrafo correcto para una configuración puede no serlo para otra. Se comprueban suplementos, referencias, revisión y efectividad. La revisión más reciente no es necesariamente la exigida por una AD que incorpora una concreta y no permite posteriores automáticamente.

AMM y EMM cubren aeronave y motor; CMM trata componente definido. IPC identifica relaciones, sin ser instrucción universal de reparación; SRM aporta datos dentro de límites. MPD y MRBR contribuyen a planificación en su función, sin equivaler al programa de cualquier operador. ALS contiene limitaciones en su alcance.
\section{Orden de trabajo y ejecución real}
Una orden identifica producto, configuración, origen, datos y revisión, acción, recursos, ejecución, resultados y cierre. Conserva referencias inequívocas. Trazabilidad significa saber qué se hizo, no solo qué estaba previsto. Un hallazgo necesita resolución por la vía pertinente.

No se escoge el valor más cómodo ni se interpola sin autorización entre instrucciones de configuraciones diferentes. Un conflicto requiere aclaración atendiendo a naturaleza y función de documentos; no se resuelve solo por el número más alto de revisión.
\section{Historial y corrección controlada}
La identidad y estado de una pieza siguen sus instalaciones y retiradas. El documento de liberación informa en su alcance; elegibilidad añade otra comprobación. Un resumen no sustituye automáticamente la evidencia original requerida.

Corregir conserva qué cambió, motivo, fuente y responsabilidad. Borrar para que la hoja cuadre puede destruir trazabilidad. La integridad de datos es relevante para seguridad: unidades o fórmulas incorrectas pueden producir falsa apariencia de margen.
\repaso{¿Figurar en IPC demuestra que una pieza está lista para instalar?}{Se comprueban estado, historial, modificaciones y elegibilidad para la configuración. La referencia de catálogo solo identifica una parte de la cuestión.}
